\documentclass{article}
\usepackage[tmargin=0.5cm,rmargin=0.5in,lmargin=0.5in,margin=0.5in,bmargin=2cm,footskip=.2in]{geometry}

\usepackage{import}

\title{\textbf{Pedago - PKE}}
\date{\textit{Un beau lundi de 2025}}
\author{Hasley William}

\usepackage{nameref}
\usepackage[ruled, vlined,linesnumbered, french, frenchkw]{algorithm2e}
\usepackage{setspace}
\usepackage{comment} % enables the use of multi-line comments (\ifx \fi) 

\SetKwInput{KwRes}{Sortie}%
\SetKwIF{Si}{SinonSi}{Sinon}{si}{alors}{sinon si}{sinon}{fin si}%
\SetKwFor{Pour}{Pour}{faire}{fin Pour}%
\SetKwInput{KwIn}{Entrée}
\SetKwInput{KwOut}{Sortie}

\usepackage[none]{hyphenat}
\usepackage[utf8]{inputenc}
\usepackage{tikz}
\usepackage{tkz-tab}
\usepackage{amsfonts} 
\usepackage{amsmath}
\usepackage{amssymb}
\usepackage[mathscr]{eucal}
\usepackage{bigints}


\usepackage[upright]{fourier} % PACKAGE QUI MODIFIE LA FONT DE BASE

\usepackage{eulervm} % Fixe les matrices mais détruit les inégalités et les cases -> A fix à la main
\usepackage{euler} %eulervm
\usepackage{mathdots}

\usepackage{stmaryrd}

\usepackage{fancyhdr}
\usepackage{graphicx}
\usepackage{fancybox}
\usepackage[most]{tcolorbox}
\usepackage[utf8]{inputenc}
\usepackage[T1]{fontenc}
\usepackage{mathtools, array, boldline}
\usetikzlibrary{automata, positioning, shapes}

\parindent 0ex


\begin{document}

\begin{algorithm}[H]
\setstretch{1.35}
\KwIn{La carte publique du receveur. Un nombre à envoyer}
\KwOut{Le message encodé}
\SetAlgoLined
\SetNoFillComment
\vspace{3mm}

Choisir des nombres aléatoirement pour chaque point de la carte\;
Ajuster ces choix afin que le total de ces nombres corresponde au message à envoyer\;

\Pour{\textsl{Chaque point de la carte}}{
	Modifier la valeur associée à ce point en encodant la somme des valeurs voisines (avec la valeur actuelle)\;
}

Envoyer la carte avec ces nombres au receveur\;

\caption{Procédure d'Encodage}
\end{algorithm}

\vspace{5em}

\begin{algorithm}[H]
\setstretch{1.35}
\KwIn{La carte privée du receveur, remplie par l'envoyeur. Un nombre à décoder}
\KwOut{Le message décodé}
\SetAlgoLined
\SetNoFillComment
\vspace{3mm}

Sommer chaque nombre sur les points spéciaux de la carte\;
Retrouver le nombre envoyé\;

\caption{Procédure de Décodage}
\end{algorithm}

\vspace{3em}

\noindent\rule{\textwidth}{1pt}

\vspace{3em}

\begin{algorithm}[H]
\setstretch{1.35}
\KwIn{La carte publique du receveur. Un nombre à envoyer}
\KwOut{Le message encodé}
\SetAlgoLined
\SetNoFillComment
\vspace{3mm}

Choisir des nombres aléatoirement pour chaque point de la carte\;
Ajuster ces choix afin que le total de ces nombres corresponde au message à envoyer\;

\Pour{\textsl{Chaque point de la carte}}{
	Modifier la valeur associée à ce point en encodant la somme des valeurs voisines (avec la valeur actuelle)\;
}

Envoyer la carte avec ces nombres au receveur\;

\caption{Procédure d'Encodage}
\end{algorithm}

\vspace{5em}

\begin{algorithm}[H]
\setstretch{1.35}
\KwIn{La carte privée du receveur, remplie par l'envoyeur. Un nombre à décoder}
\KwOut{Le message décodé}
\SetAlgoLined
\SetNoFillComment
\vspace{3mm}

Sommer chaque nombre sur les points spéciaux de la carte\;
Retrouver le nombre envoyé\;

\caption{Procédure de Décodage}
\end{algorithm}


\end{document}
